% Figure: centroid of a quarter disc computed in polar coordinates. The
% region is the quarter disc of radius R in the first quadrant. The polar
% area element r dr dtheta and the moment element weight the calculation;
% the centroid sits at (4R/3pi, 4R/3pi), marked.
\begin{figure}[htbp]
\centering
\begin{tikzpicture}
\begin{axis}[
  width=8.6cm, height=8.0cm,
  xlabel={$x$}, ylabel={$y$},
  xmin=-0.15, xmax=1.25, ymin=-0.15, ymax=1.25,
  axis lines=left, font=\small,
  axis equal image,
]
% quarter disc R = 1, shaded
\addplot[draw=none, fill=engblue!15, domain=0:90, samples=80]
  ({cos(x)}, {sin(x)}) \closedcycle;
% arc
\addplot[engblue, very thick, domain=0:90, samples=80]
  ({cos(x)}, {sin(x)});
% straight edges
\addplot[engblue, very thick] coordinates {(0,0) (1,0)};
\addplot[engblue, very thick] coordinates {(0,0) (0,1)};
% polar area element sketch: a small wedge near (r,theta)=(0.6,42 deg)
% corner coordinates precomputed to avoid macros inside coordinates{}
\addplot[engamber, thick, fill=engamber!25, samples=2]
  coordinates {(0.4334,0.3386) (0.5516,0.4309)
               (0.4863,0.5035) (0.3820,0.3956)}
  \closedcycle;
\node[font=\scriptsize, engamber, anchor=west]
  at (axis cs:0.72,0.55) {$r\,\dd r\,\dd\theta$};
% centroid at (4/(3pi), 4/(3pi)) = (0.4244, 0.4244)
\addplot[engred, only marks, mark=*, mark size=2pt]
  coordinates {(0.4244,0.4244)};
\node[font=\scriptsize, engred, anchor=south west]
  at (axis cs:0.44,0.44) {$(\bar{x},\bar{y})$};
\end{axis}
\end{tikzpicture}
\caption[Centroid of a quarter disc in polar coordinates]{The centroid
of a quarter disc of radius $R$, computed cleanly in polar coordinates
where the region is the coordinate box $0 \leq r \leq R$, $0 \leq \theta
\leq \pi/2$. The polar area element $\dd A = r\,\dd r\,\dd\theta$ (amber
wedge) carries the Jacobian factor $r$ that a Cartesian description would
have to recover with an awkward $\sqrt{R^{2} - x^{2}}$ boundary. The
first moment $\iint x\,\dd A = \int_{0}^{\pi/2}\!\int_{0}^{R}
(r\cos\theta)\,r\,\dd r\,\dd\theta = R^{3}/3$ divided by the area
$\pi R^{2}/4$ places the centroid at $\bar{x} = \bar{y} = 4R/3\pi
\approx 0.424\,R$ (red point). Choosing the coordinate system that
straightens the boundary turns a two-page Cartesian integral into two
elementary polar integrations.}
\label{fig:vol02ch06:polar-centroid}
\end{figure}
